\section{From ordinary end data to every coefficient-algebra point}
\label{sec:unramified-coefficient-extension-v49}

This chapter reconstructs the original curve and specialization arguments
between ordinary end data and the joint unramified-coefficient automorphy
assertion.  The conclusions are the literal \texttt{FinCrit},
\texttt{ConstOfChar}, \texttt{ordPt\_of\_gd},
\texttt{finA\_of\_end}, \texttt{exists\_ordPt\_of\_point},
\texttt{matIrr\_open\_gd}, and the same-marked-element implication
\[
 \mathrm{GdEnd}\Rightarrow\mathrm{AutomorphyUnramifiedCoefficients}.
\]
The common ordinary end datum remains an explicit input to the last
implication.  In particular, the chapter does not assert the existence of
its finite-characteristic or automorphic end conditions.

Here are the objects and the exact conventions.  Fix the original
\(D:\mathrm{Dat}(\Omega,K,c,e,\psi)\), put
\(\mathcal O=c.O\), and write
\[
 G=G_{\mathbf Q},\qquad H=\ker\eta=G_K,\qquad
 \Psi=\epsilon^e\psi,\qquad (w_0,w_1,w_2)=(1,0,-1).
 \tag{UCE.1}
\]
All algebraic closures, embeddings and the later bridge
\(\iota:\overline{\mathbf Q}_\ell\simeq\mathbf C\) are the supplied
ones.  The group is
\(\mathcal G_3=\mathrm{GL}_3\rtimes\{1,-1\}\), with
\(\theta(M)=JM^{-t}J\), where \(J\) is the antidiagonal permutation
matrix.  Put \(T(t)=\operatorname{diag}(t,1,t^{-1})\) and
\(r_0(g)=(T(\Psi(g)),\eta(g))\).  Entries have weights
\(w_i-w_j\).  The fixed complex conjugation is denoted by \(\sigma\).
A marked element satisfies
\[
 \sigma^2=1,\quad \eta(\sigma)=-1,\quad \Psi(\sigma)=1,
 \quad\eta(g_1)=1,\quad \Psi(g_1)^2-1\in\mathcal O^\times.
 \tag{UCE.2}
\]
Normalization means that the matrix of \(r(\sigma)\) is one and the
\((0,1)\) and \((1,0)\) entries of
\(r(g_1)-r(\sigma g_1^{-1}\sigma)\) vanish.  We always use the actual
local ideal \(I_t\) and the actual normalized universal ring \(S\)
of OCH7--OCH9.

The original \(GD\) assumptions remain unchanged: \(\ell\ne2\),
\(e\ge1\), even nonzero \(N\), prime-to-\(\ell\) order of
\(\psi\), \((\operatorname{comp}_\ell\psi)^2=1\), its prescribed
parity, imaginary quadratic \(K\),
\(\ell\nmid |\operatorname{disc}K|h_K\),
\((|\operatorname{disc}K|,N)=1\), the original residual kernel
noncontainment, and \(\ell-1\nmid2e\).  No lower bound on the prime
is added.  The local and finite-characteristic lemmas include prime
three in their stated scopes; the last \(GD\) condition itself excludes
it in the final arithmetic application.

For a valuation ring \(B\subset\Omega\) over \(\ell\), an
\emph{ordinary point} over a complete local algebra \(A\) has the
five original fields: the prescribed residual representation and
component character; continuity modulo every power of \(\mathfrak m_A\);
the normalization just described; triviality on
\(I_{B'}\cap\ker\eta\cap\ker\Psi\) at every place away from
\(\ell\); and, at \(B\), a basis congruent to the identity in which
the representation on the kernel of \(\eta\) in the decomposition
group is upper triangular, with inertia diagonal
\((\Psi,1,\Psi^{-1})\).  This is exactly \texttt{OrdPt}.

The already reconstructed inputs used below are OCH4--OCH9 and the
local part of OCH15, with their named period and cohomology foundations.
The additional background is elementary complete local algebra,
valuation rings, finite extensions of complete discretely valued
fields, and the usual Galois decomposition/inertia dictionary.  The
proofs below give the curve detection, finite generation, ordered
matrix arguments and extension of arbitrary algebra maps explicitly.

\subsection{Ordinary points over the original coefficient rings}

% BEGIN UCE1
\paragraph{UCE1. The actual inertia separator and marked elements.}
\label{claim:UCE1}
For every valuation ring \(B\) of \(\Omega\) above \(\ell\), there is
\(x\in I_B\) such that
\(\eta(x)=1\) and \(\Psi(x)^2-1\in\mathcal O^\times\).
In particular, taking \(g_1=x\) gives (UCE.2).  This conclusion alone
does not assign the end conditions to that choice.

\emph{Proof.}
The cyclic group \(\mathbf F_\ell^\times\) has an element \(a\)
with \(a^{2e}\ne1\), since \(\ell-1\nmid2e\).
Inertia surjects onto the cyclotomic character modulo \(\ell\):
\(\mathbf Q_\ell(\mu_\ell)/\mathbf Q_\ell\) is totally ramified.
Choose an inertia element with cyclotomic reduction \(a\).
The character \(\eta\) is unramified at \(\ell\), and the square
of the finite inertia part of \(\Psi\) is one by the original
\(\operatorname{comp}_\ell\psi\) condition.  Consequently
\(\overline\Psi(x)^2=a^{2e}\ne1\).  This is precisely the asserted
unit condition.  Changing the prolongation of the place conjugates
inertia; both \(\eta\) and \(\Psi\) are characters and retain their
values.  The other marked identities are the supplied complex
conjugation and parity identities.  This proves the assertion on the
given algebraic closure and for every original \(B\).
% END UCE1

% BEGIN UCE2
\paragraph{UCE2. Ordered flags specialize after a finite field extension.}
\label{claim:UCE2}
Let \(R\) be a domain, let matrices \(M_s\) and prescribed diagonal
entries \(d_{s,i}\) be defined over \(R\), and suppose that over a
field containing \(\operatorname{Frac}R\) they admit one ordered
upper triangular flag, with those diagonals for every index in the
specified inertia subset.  For every map \(R\to k\) to a field,
the specialized matrices admit the same ordered triangular and
diagonal conditions over a finite extension of \(k\).
No restriction is imposed on the characteristic of \(k\).

\emph{Proof.}
Let \(\mathfrak p\) be the kernel of the map.  Choose a valuation
ring in the generic field dominating \(R_{\mathfrak p}\).
To recall the existence argument, order the local subrings of the
field dominating \(R_{\mathfrak p}\) by domination.  A union of a
chain is again such a ring.  A maximal one is a valuation ring:
if neither \(z\) nor \(z^{-1}\) belongs to it, adjoining one and
localizing at a prime above its maximal ideal gives a larger
dominating local ring, a contradiction.  Denote it by \(V\).

Every flag in the fraction field of \(V\) has a basis matrix
\(k b\), where \(k\in\mathrm{GL}_n(V)\) and \(b\) is upper
triangular over the field.  Indeed, in the first column choose an
entry dividing all its entries; such an entry exists because the
finite set of principal ideals is totally ordered.  Scale the
column to a primitive vector, carry it to the first basis vector by
invertible row operations over \(V\), and continue in the quotient.
Induction gives the assertion.  Changing the generic flag basis by
\(b\) preserves upper triangularity and all ordered diagonal entries.
In the basis \(k\), the matrices are integral because the original
matrices lie in \(V\).  Their asserted entries therefore remain zero,
or equal to the prescribed diagonals, on reduction to \(\kappa(V)\).

Domination identifies the kernel of \(R\to\kappa(V)\) with
\(\mathfrak p\).  The two fields \(k\) and \(\kappa(V)\) have a
common extension over \(\operatorname{Frac}(R/\mathfrak p)\): take
a prime quotient of their nonzero tensor product and then its
fraction field.  This gives a flag over some extension of \(k\).
Finally consider the affine scheme of invertible basis matrices
with the required triangular and diagonal entries.  Its equations
are polynomial, after adding an inverse for the determinant.
Although their indices may be infinite, their ideal in the
finite-variable polynomial ring is finitely generated.  It is proper
because of the flag just constructed.  A maximal ideal containing it
has residue field finite over \(k\), by the finite-type field lemma.
That point supplies all the original equations simultaneously and
proves the required finite-extension assertion.
% END UCE2

% BEGIN UCE3
\paragraph{UCE3. The full local ordered-flag assertion in every characteristic.}
\label{claim:UCE3}
For every original local datum \(D_\ell:\mathrm{Dat3}\), its
universal ring \(U\), every field \(k\) over \(\mathcal O\),
and every actual representation obtained from a map
\(U/I_t\to k\), there is a finite extension \(k'/k\) and a
basis over \(k'\) giving the full original \texttt{TriIn} condition.
Thus the original \texttt{lemma2c} is supplied, including targets
of characteristic \(\ell\).

\emph{Proof.}
OCH7 gives the domain \(R=U/I_t\), its injective characteristic-zero
formal chart, and its base point.  In every finite truncation of that
chart, OCH5 makes the representation crystalline and OCH4 makes it
ordinary.  Its diagonal infinitesimal characters, after division by
the fixed characters, are unramified: the weight-zero crystalline
\(H^1_f\) calculation used in OCH4 and OCH15 has precisely that
conclusion.  Induction on the Artin length therefore keeps the
\emph{full} inertia diagonals equal to
\(\lambda_i(s)=\Psi(s)^{w_i}\).
Consequently the following polynomial identities hold in every
truncation and hence in the injective chart and in \(R\):
\[
 \det(X-r(s))=\prod_{i=0}^2(X-\lambda_i(s)),\qquad
 \prod_{i=0}^2\bigl(r(s_i)-\lambda_i(s_i)I\bigr)=0
       \quad(s,s_i\in I_\ell),
 \tag{UCE.3}
\]
where the second product has the displayed order.

Over \(\operatorname{Frac}R\), choose an inertia element \(u\)
in the kernel of the finite character part whose cyclotomic value
has infinite order.  Such an element exists because an open subgroup
of inertia still has infinite cyclotomic image.  Since \(e\ge1\),
the three values \(\lambda_i(u)\) are distinct in characteristic
zero.  The first identity in (UCE.3) diagonalizes \(r(u)\) in that
order.  In that basis apply the second identity successively to
\((s,u,u),(u,s,u),(u,u,s)\).  Division by the nonzero eigenvalue
differences shows that \(r(s)\) is upper triangular and its
\(i\)-th diagonal is \(\lambda_i(s)\).  For example, the first
product kills the entries below the first diagonal position and
identifies its diagonal entry; the last product does the same for
the last row, and the middle product identifies the middle diagonal.

If \(g\) is in the kernel of \(\eta\) in the decomposition
group, then \(gug^{-1}\) is in inertia and has the same ordered
character values.  The equation
\(r(gug^{-1})r(g)=r(g)r(u)\), solved from the bottom row upward,
forces \(r(g)\) to be upper triangular.  At each off-diagonal
step the coefficient is one of the same nonzero eigenvalue
differences.  We have obtained the full ordered generic flag,
with all decomposition and inertia equations, over
\(\operatorname{Frac}R\).  Apply UCE2 to the given map \(R\to k\).
The multiplier is unchanged throughout.  Its output is exactly the
original finite-extension \texttt{TriIn} assertion.
% END UCE3

% BEGIN UCE4
\paragraph{UCE4. Descent of the ordered flag to the actual integral basis.}
\label{claim:UCE4}
Let \(A\) be a complete local domain over \(\mathcal O\), and let
\(r\) be a residual lift of the actual local rank-three datum over
\(A\).  Suppose its fraction-field representation has the ordered
\texttt{TriIn} flag after a field extension.  If inertia contains
\(x\) with \(\Psi(x)^2-1\in\mathcal O^\times\), then \(r\)
satisfies the original integral \texttt{CondT}: its basis can be
chosen over \(A\), congruent to the identity.

\emph{Proof.}
Put \(d_i=\Psi(x)^{w_i}\).  Every \(d_i-d_j\), \(i\ne j\),
is a unit: each differs by a unit factor from \(\Psi(x)-1\),
\(\Psi(x)+1\), or \(\Psi(x)^2-1\).
The generic ordered flag gives
\(\det(X-r(x))=\prod_i(X-d_i)\); injectivity into the field
puts the same identity in \(A[X]\).  Form the spectral projectors
\[
 P_i=\prod_{j\ne i}\frac{r(x)-d_jI}{d_i-d_j}\quad\in M_3(A).
 \tag{UCE.4}
\]
Cayley--Hamilton makes them orthogonal idempotents of sum one.
Their reductions are the three coordinate projectors.  Thus
\(P_i e_i\) has \(i\)-th coordinate a unit; divide by that unit.
These three vectors form a matrix \(h\equiv I\pmod{\mathfrak m_A}\)
and \(h^{-1}r(x)h=\operatorname{diag}(d_i)\).

Let \(k\) be a generic ordered-flag matrix.  The equation
\[
       (k^{-1}r(x)k)(k^{-1}h)=(k^{-1}h)\operatorname{diag}(d_i)
\]
forces \(k^{-1}h\) to be upper triangular, by solving its entries
from the bottom row upward.  The different \(d_i\) have already
been matched in their prescribed order.  Conjugation by an upper
triangular matrix preserves upper triangularity and its diagonal
characters.  The required equations therefore hold in the \(h\)
basis over the generic field, and then over \(A\) by injectivity.
This proves every decomposition-group and inertia clause of
\texttt{CondT}, with the actual integral near-identity matrix.
% END UCE4

% BEGIN UCE5
\paragraph{UCE5. Every original local domain specialization is an ordinary point.}
\label{claim:UCE5}
Let \(D,g_1,U,S\) be original data with \(g_1\) marked, and let
\(B\subset\Omega\) be the actual valuation ring defined by
\(x\in B\Longleftrightarrow\|D.j(x)\|\le1\).
For every complete local domain \(A\) over \(\mathcal O\) and
every local algebra map \(f:S\to A\), the actual push
\(f_*r_S\) is an \(\mathrm{OrdPt}\) for the same
\(D,g_1,B\).  This is the full original \texttt{ordPt\_of\_gd}.

\emph{Proof.}
A local map sends \(\mathfrak m_S^a\) into
\(\mathfrak m_A^a\) for every \(a\).  Hence the residual matrix
identities and all locally constant reductions of \(r_S\) push
to the exact residual and adic-continuity fields.  The normalization
equations also push literally.
At a place away from \(\ell\), an element with
\(\eta(s)=\Psi(s)=1\) belongs to the original killed-inertia
set, so its universal image, and therefore its image under \(f\),
is one.  This proves the original all-place away condition.

The specified local map has image the decomposition group of \(B\)
and takes its local inertia onto \(I_B\); this is the fixed
embedding dictionary in OCH9.  The local factorization of \(S\)
kills the same \(I_t\).  Apply UCE3 over \(\operatorname{Frac}A\),
then UCE1 and UCE4.  The resulting integral basis gives exactly
the last field of \(\mathrm{OrdPt}\).  This works even when
\(\mathcal O\to A\) kills its maximal ideal, and uses neither
\(\mathrm{FinA}\) nor an automorphy hypothesis.
% END UCE5

\subsection{The actual weighted ring and its characteristic curves}

% BEGIN UCE6
\paragraph{UCE6. The integral weighted presentation by actual entries.}
\label{claim:UCE6}
For every original normalized ring \(S\), there is a finite
weighted presentation
\[
 P=\mathcal O[[X_1,\ldots,X_d]]\twoheadrightarrow S,
 \qquad X_i\longmapsto z_i,
 \tag{UCE.5}
\]
where each \(z_i\in\mathfrak m_S\) is an actual matrix-entry
deviation from \(r_0\), its weight \(a_i\) is the corresponding
\(w_b-w_c\), and the kernel is homogeneous under every formal
weight projection.  Every universal matrix entry belongs to its
asserted homogeneous piece.  These are the literal \texttt{WPres}
properties needed in both original coefficient arguments.

\emph{Proof.}
The entry deviations span the relative cotangent space.  Otherwise
a nonzero functional on that space vanishing on all the deviations
would give two different maps to the residue dual-number algebra
with identical pushed representations, contradicting the uniqueness
in the actual universal property.  Choose finitely many deviations
spanning that finite-dimensional space.  Successive approximation
and completeness, beginning with Nakayama's lemma, give the
surjection (UCE.5).

Here is a proof of integral homogeneity that does not divide by
differences of torus eigenvalues.  Over \(S[[T]]\), conjugate the
universal representation by \(T(1+T)\).  This preserves the
residual representation, killed inertia and normalization; OCH8
preserves its literal local \(I_t\)-condition.  Universality gives
a map \(S\to S[[T]]\) that multiplies every entry of weight \(j\)
by \((1+T)^j\), and specializes at \(T=0\) to the identity.
For \(F\in\ker(P\to S)\), it follows that
\(F((1+T)^{a_i}z_i)=0\).
Modulo \(\mathfrak m_S^b\), only finitely many monomials in the
\(z_i\) survive.  Replace \(1+T\) by a Laurent indeterminate
\(Z\).  The substitution
\[
 (S/\mathfrak m_S^b)[Z,Z^{-1}]
       \longrightarrow(S/\mathfrak m_S^b)[[T]],\qquad Z\mapsto1+T
\]
is injective: after multiplication by a power of \(Z\), this is
the invertible integral change of variable in a polynomial.
Each coefficient of \(Z^j\) is therefore zero.  Letting \(b\)
increase shows that every weight projection \(F_j\) is in the
original kernel.  The same argument applied to a lift of any
universal entry shows that all its other weight projections vanish.
Thus the pieces \(S_j\), defined as images of homogeneous formal
series, have independent finite sums, multiply in the prescribed
weights, and contain the actual entries with their original weights.
The chosen deviations also lie in the span of entry monomials of
their weight, because their base constants vanish off weight zero.
% END UCE6

% BEGIN UCE7
\paragraph{UCE7. Weight zero and every fixed weight have finite monomial presentations.}
\label{claim:UCE7}
For any finite weighted presentation as in (UCE.5), with an integral
homogeneous kernel and variables in the maximal ideal, \(S_0\) is
a complete Noetherian local \(\mathcal O\)-algebra with the same
residue field.  Each \(S_j\) is a finite \(S_0\)-module.  The
maximal ideal of \(S_0\) is generated by \(\mathfrak m_{\mathcal O}\)
and finitely many positive-degree weight-zero monomials.  Moreover
\[
               S_0\cap\pi S=\pi S_0
 \tag{UCE.6}
\]
for a uniformizer \(\pi\) of \(\mathcal O\).

\emph{Proof.}
The nonzero minimal exponent vectors of weight zero under
componentwise comparison are finite by the ascending-chain property
of monomial ideals in a finite-variable polynomial ring.  Every
nonzero weight-zero exponent vector dominates one of them; subtract
it and repeat.  Thus they generate the weight-zero monoid.  For
a fixed weight \(j\), the minimal vectors of that weight are
again finite, and every vector of weight \(j\) is a minimal one
plus a weight-zero vector.  Grouping formal monomials by these
choices shows that the weight-zero power-series ring is a quotient
of a finite-variable power-series ring over \(\mathcal O\), and
that each fixed-weight series module is finite over it.  The
grouping converges: the lengths of the chosen factorizations tend
to infinity with total degree, since the finite generator degrees
are bounded.

The homogeneous kernel in (UCE.5) gives the same statements after
quotienting.  The kernels in these finite-module presentations are
closed, so their images are complete.  A weight-zero element is
a constant from \(\mathcal O\) plus an element generated by the
positive-degree weight-zero monomials.  If the residue of the
constant is nonzero, its inverse is the convergent geometric series,
still of weight zero.  This identifies the local ring, its residue
field and its asserted maximal ideal.

Finally, if \(s\in S_0\) and \(s=\pi y\) in \(S\), choose
a homogeneous weight-zero lift \(F\) of \(s\) and an arbitrary
lift \(Y\) of \(y\).  Project \(F-\pi Y\), which lies in
the homogeneous kernel, to weight zero.  It gives
\(s=\pi\overline{Y_0}\) in \(S_0\).  This proves (UCE.6),
including rings with \(\pi\)-torsion.
% END UCE7

% BEGIN UCE8
\paragraph{UCE8. Detection of a surviving element by a characteristic-prime DVR curve.}
\label{claim:UCE8}
Let \(R\) be a complete Noetherian local ring with finite residue
field of characteristic \(p\).  Let \(Q\) be a prime containing
\(p\), and let \(b\in\mathfrak m_R\setminus Q\).  There is
a local map \(R\to A\) to a complete discrete valuation ring of
characteristic \(p\), with finite residue field, whose kernel
contains \(Q\) and does not contain \(b\).

\emph{Proof.}
Choose a prime \(\mathfrak p\supset Q\), maximal among primes
not containing \(b\), using the ascending-chain condition.
Then \(D=R/\mathfrak p\) is a complete local domain and the
nonzero image of \(b\) lies in its maximal ideal.  Its dimension
is one.  Here is the dimension argument.  Every nonzero prime of
\(D\) contains \(b\), by maximality.  For a nonzero
\(x\in\mathfrak m_D\), a prime minimal over \((x)\) has
height one by the principal ideal theorem.  It contains \(b\),
and hence is also minimal over \((b)\), since a strict intermediate
nonzero prime would give height at least two.  Consequently
\(\mathfrak m_D\) is contained in the union of the finitely many
primes minimal over \((b)\).  Finite prime avoidance puts it in
one of them.  It has height one, as asserted.

For completeness, the required discrete valuation ring can be
constructed directly in this equal-characteristic situation.
Let \(k_D\) be the finite residue field of \(D\), of order \(q\).
The roots of \(X^q-X\) lift uniquely by completeness and its unit
derivative \(-1\).  In characteristic \(p\) they are closed
under addition, multiplication and inverse, and give a coefficient
field \(k_D\subset D\).  Choose any nonzero
\(t\in\mathfrak m_D\).  Dimension one makes \((t)\)
\(\mathfrak m_D\)-primary, so the \(t\)-adic and maximal-ideal
topologies are equivalent.  The map \(k_D[[T]]\to D\),
\(T\mapsto t\), is injective: a nonzero series is a power of
\(T\) times a unit, and neither factor becomes zero.  Lifting a
finite basis of \(D/tD\) and successively subtracting its constant
terms proves that \(D\) is finite over \(k_D[[t]]\).

Thus \(\operatorname{Frac}D\) is a finite extension of
\(k_D((t))\).  The extension of its complete discrete valuation
has a complete DVR of integers \(A\), finite over \(k_D[[t]]\)
with finite residue field.  Every element of \(D\) is integral
over \(k_D[[t]]\), so belongs to \(A\).  The inclusion is local:
a power of \(\mathfrak m_D\) lies in \(tD\), whereas units
stay units.  Its injectivity keeps \(b\) nonzero.  Composing the
quotient map with this inclusion proves the assertion.  The same
argument applies when the specified prime contains a uniformizer
of a coefficient DVR: the quotient then has the residue
characteristic, which is the case used below.
% END UCE8

% BEGIN UCE9
\paragraph{UCE9. The complete original finite weight-zero criterion.}
\label{claim:UCE9}
Let \(\mathcal O\) be any complete DVR with finite residue field,
let \(R\) be an original complete Noetherian local
\(\mathcal O\)-algebra with the same residue field, and let
\(R\) have a finite integral weighted presentation with all
variables in its maximal ideal.  Suppose that for every local map
\(f:R\to A\) to a complete DVR with finite residue field, for
which \(f(\mathfrak m_{\mathcal O})=0\), all weight-zero
variables vanish and either all negative-weight variables or all
positive-weight variables vanish.  Then \(R_0\) is finite over
\(\mathcal O\).  This is the literal \texttt{FinCrit}, with all
its original test rings and arbitrary integer variable weights.

\emph{Proof.}
Take one of UCE7's finitely many positive-degree weight-zero
monomial generators \(b\).  If its image in \(R\) were not
nilpotent modulo \(\pi R\), a prime \(Q\supset\pi R\)
would avoid it.  UCE8 supplies an original test map with
\(f(b)\ne0\).  In a nonempty weight-zero monomial, either a
weight-zero variable occurs, or both a negative and a positive
weight occur.  The assumed alternative therefore kills at least
one of its factors.  This contradiction proves
\(b^m\in\pi R\) for some \(m\).
Equation (UCE.6) puts \(b^m\) in \(\pi R_0\).

All generators of the maximal ideal of \(R_0\) are consequently
in \(\sqrt{\pi R_0}\).  Since there are finitely many, some
power of that maximal ideal lies in \(\pi R_0\).  The quotient
\(R_0/\pi R_0\) is therefore Artinian with finite residue field,
and is finite-dimensional over \(\mathcal O/\pi\).
The two adic topologies on \(R_0\) are equivalent.  Lift a finite
set of generators of this quotient.  At each successive
\(\pi\)-power subtract their \(\mathcal O\)-linear combination;
completeness sums the resulting coefficients in \(\mathcal O\).
The lifts generate \(R_0\) as an \(\mathcal O\)-module.
This proves finiteness, without assuming that \(R\) is reduced,
flat over \(\mathcal O\), or characteristic zero.
% END UCE9

% BEGIN UCE10
\paragraph{UCE10. Finite-index characteristic polynomials determine the constants everywhere.}
\label{claim:UCE10}
Let \(k\) be a finite field, let \(A\) be a local domain over
\(k\), and let a representation over \(A\) reduce to the supplied
diagonal character representation over \(k\).  If their
characteristic polynomials agree on a finite-index subgroup, then
they agree at every group element.  For a disconnected rank-three
representation the assertion applies on \(H=\ker\eta\), exactly
as in the original \texttt{ConstOfChar} contract.

\emph{Proof.}
For each \(g\), some positive power \(g^a\) lies in the
given finite-index subgroup.  Every eigenvalue of \(r(g)^a\)
is algebraic over \(k\), so every eigenvalue of \(r(g)\) is
algebraic over \(k\).  Hence all coefficients of its characteristic
polynomial are algebraic over \(k\).
An element \(y\in A\) algebraic over \(k\) and congruent to
\(c\in k\) is equal to \(c\): the domain \(k[y]\) is a
finite-dimensional field over \(k\); if \(y-c\ne0\), its
inverse belongs to \(k[y]\subset A\), contrary to membership
in the maximal ideal.  Applying this argument to every coefficient
gives the asserted equality with the residual polynomial.
On \(H\), the matrix of \(r(g^a)\) is the ordinary power
of the matrix of \(r(g)\), so the same proof applies literally.
It uses neither continuity nor a prime-to-characteristic index.
% END UCE10

% BEGIN UCE11
\paragraph{UCE11. Constant characteristic polynomials force the original triangular alternative.}
\label{claim:UCE11}
The full original \texttt{ConstOfChar} holds.  Namely, over a local
domain \(A\) over a finite field \(k\) of characteristic different
from two, take the supplied characters \(\eta,\Psi\), marked
\(\sigma,g_1\), normalized residual lift \(r\), and any
finite-index subgroup on which the characteristic polynomials are
those of \(r_0\).  Then all diagonal entries of \(r\) are the
given base entries, and \(r\) is either upper triangular on \(H\)
or lower triangular on \(H\).  No characteristic-three exception
is imposed.

\emph{Proof.}
UCE10 first gives the base characteristic polynomial at every
element of \(H\).  Normalization and the disconnected transpose
identity give, over \(A\),
\[
 r(g_1)-r(\sigma g_1^{-1}\sigma)
       =\zeta\operatorname{diag}(1,0,-1).
 \tag{UCE.7}
\]
Indeed the difference is minus its antidiagonal transpose, and its
two prescribed adjacent entries are zero.  The entries fixed by
that transpose are zero because two is invertible.  The residual
value of \(\zeta\) is \(\Psi(g_1)-\Psi(g_1)^{-1}\ne0\),
so \(\zeta\) is a unit.

Work over \(F=\operatorname{Frac}A\).  Every invariant subspace
is stable under (UCE.7), and hence is a coordinate subspace by
its three spectral projectors.  A zero \((i,j)\) entry forced
by such a subspace also forces the zero \((2-j,2-i)\) entry,
because
\(r(\sigma g^{-1}\sigma)=Jr(g)^tJ\) on \(H\).
The six proper nonempty coordinate subsets therefore give exactly
the following alternatives: the whole representation is irreducible;
it is upper triangular; it is lower triangular; or all four
odd-weight entries vanish and its outer two-dimensional block is
irreducible.  If that outer block were reducible, its two distinct
marked projectors would give a coordinate line and one of the
triangular alternatives.

We rule out both irreducible alternatives by an explicit trace
Gram argument.  For an irreducible block of size \(n=3\), or
the outer block of size \(n=2\), its marked diagonal has distinct
eigenvalues.  The coordinate matrix projectors belong to the
\(F\)-algebra spanned by its group matrices.  The directed graph
whose edges are nonzero matrix entries is strongly connected:
the coordinate span reachable from any vertex would otherwise be
a proper invariant subspace.  Multiplying a matrix on both sides
by coordinate projectors gives its matrix unit times a nonzero
scalar.  Multiplying along paths therefore supplies every matrix
unit.  The span of the group matrices is thus all \(M_n(F)\).

Choose \(n^2\) group matrices \(M_i\) forming an \(F\)-basis.
The matrix
\[
                     G_{ij}=\operatorname{tr}(M_iM_j)
 \tag{UCE.8}
\]
is nonsingular because the trace pairing on \(M_n(F)\) is
nondegenerate, in every characteristic.  For the full block its
entries lie in \(k\), by UCE10.  For the outer block, first note
that the middle entry is an eigenvalue of the full block matrix;
it is algebraic over \(k\) and congruent to one, so is one by
UCE10's elementary field argument.  The outer trace is then the
full trace minus one, and again (UCE.8) has entries in \(k\).
Its nonzero determinant is a nonzero constant, and hence remains
nonzero in the residue field of \(A\).  But every reduced \(M_i\)
is diagonal, so the reduced matrix (UCE.8) factors through the
\(n\)-dimensional space of diagonal matrices.  Its rank is at
most \(n<n^2\), a contradiction.  This argument does not divide
by \(n\) or \(n!\).

We have proved the triangular alternative.  Each diagonal entry
is now a root of the constant characteristic polynomial; it is
algebraic over \(k\) and has its prescribed residual value.
UCE10's field argument identifies it with that value.  Finally,
for \(g\notin H\), the element \(g\sigma\) belongs to
\(H\) and has the same matrix as \(g\), because the matrix
of \(r(\sigma)\) is one.  The base matrices also agree there.
This proves all the original diagonal assertions on the full group.
% END UCE11

% BEGIN UCE12
\paragraph{UCE12. The literal finite-end condition gives finite weight zero.}
\label{claim:UCE12}
Fix the original \(D,g_1,B\), with \(g_1\) marked and \(B\)
the valuation ring supplied by the local embedding.  Suppose the
original \(\mathrm{FinEnd}\) holds: there is one finite-index
subgroup \(V\subset G\) such that every \(\mathrm{OrdPt}\)
over every complete DVR \(A\) with finite residue field and
\(\mathfrak m_{\mathcal O}A=0\) has the base characteristic
polynomials at all \(g\in V\cap H\).
Then \(\mathrm{FinA}(D,g_1)\) holds for every original local
universal ring and every original normalized global ring.

\emph{Proof.}
Take any such \(U,S\), and use UCE6's weighted presentation by
actual entry deviations.  Every local map from \(S\) to a test
DVR in the hypothesis is an actual ordinary point by UCE5.
The single given subgroup \(V\) therefore supplies the
characteristic-polynomial hypothesis of UCE11.  Since the map
kills \(\mathfrak m_{\mathcal O}\), its coefficient action
factors through the finite residue field, exactly as required there.
UCE11 kills all diagonal deviations and either all positive- or
all negative-weight entry deviations.  Thus every one of the
test maps in UCE9 satisfies its stated alternative.  UCE9 proves
that \(S_0\) is finite over \(\mathcal O\).

Every product of universal entries of total weight zero lies in
\(S_0\).  The original weight-zero entry subalgebra, being an
\(\mathcal O\)-submodule of this finite module, is finite because
\(\mathcal O\) is Noetherian.  This is precisely the defining
conclusion of \(\mathrm{FinA}\).  Neither the test subgroup nor
the marked element was changed or made dependent on \(U,S\).
This proves the full original \texttt{finA\_of\_end} implication.
% END UCE12

\subsection{Arbitrary entry-algebra points and finite coefficient realizations}

% BEGIN UCE13
\paragraph{UCE13. Under finite weight zero, the actual entry algebra is the graded algebra.}
\label{claim:UCE13}
Assume \(\mathrm{FinA}(D,g_1)\), and fix any original \(U,S\).
Let \(E_S\) be the actual entry subalgebra of \(S\), let
\(B_0\) be its original weight-zero entry subalgebra, and let
\(S_j\) be UCE6's homogeneous pieces.  Then
\[
                 B_0=S_0,\qquad E_S=\bigoplus_{j\in\mathbf Z}S_j.
 \tag{UCE.9}
\]
Each \(S_j\) is finite over \(\mathcal O\).  Its finite-module
topology agrees with its topology as a subspace of \(S\).

\emph{Proof.}
The inclusion \(B_0\subset S_0\) follows from the entry weights.
Every homogeneous weight-zero formal series is approximated, to
arbitrary maximal-ideal order, by polynomials in the entry deviations
of weight zero.  Each such polynomial belongs to \(B_0\).
By \(\mathrm{FinA}\), \(B_0\) is a finite \(\mathcal O\)-module.
It is the continuous image of a finite power of the compact ring
\(\mathcal O\) in the Hausdorff ring \(S\), and is therefore
closed.  This proves \(S_0\subset B_0\).

UCE7 expresses each \(S_j\) as a finite sum of monomials in the
chosen deviations with coefficients in \(S_0\).  Those monomials
belong to \(E_S\), so \(\bigoplus_jS_j\subset E_S\).
The reverse inclusion holds because each generating universal
entry is homogeneous.  This also shows that
\(E_S=S_0[z_1,\ldots,z_d]\), with the actual deviations.
Finiteness of each \(S_j\) follows from UCE7 and the established
finiteness of \(S_0\).  A finite free \(\mathcal O\)-module
surjects onto \(S_j\); its quotient is compact and its injection
into \(S\) is continuous and Hausdorff.  The resulting continuous
bijection is a homeomorphism.  This proves the final assertion,
including all torsion and nonreduced cases.
% END UCE13

% BEGIN UCE14
\paragraph{UCE14. Every entry-algebra point has a strictly small torus representative.}
\label{claim:UCE14}
Under the hypotheses of UCE13, let
\(\xi:E_S\to\overline{\mathbf Q}_\ell\) be any
\(\mathcal O\)-algebra map.  There is a single
\(t\in\overline{\mathbf Q}_\ell^\times\) such that the graded
map
\[
                  \xi_t(s)=t^j\xi(s)\quad(s\in S_j)
 \tag{UCE.10}
\]
takes every chosen deviation \(z_i\) to an element of strictly
positive valuation.  No continuity or localness of \(\xi\)
is assumed.

\emph{Proof.}
The finite local algebra \(S_0\) is integral over \(\mathcal O\).
Hence all its images under \(\xi\) are integral.  Moreover every
element of \(\mathfrak m_{S_0}\) has a power in \(\pi S_0\),
because the finite local algebra \(S_0/\pi S_0\) is Artinian.
Such images have strictly positive valuation.  In particular every
positive-degree weight-zero monomial in the \(z_i\) has strictly
positive valuation unless its image is zero.

For indices with \(\xi(z_i)\ne0\), write
\(b_i=v_\ell(\xi(z_i))\in\mathbf Q\).  We seek a rational
\(a\) with \(b_i+a_i a>0\) for all these indices.  Weight-zero
indices already satisfy the inequality.  If \(a_i>0>a_j\), the
balanced monomial
\(z_i^{-a_j/d}z_j^{a_i/d}\), where \(d=\gcd(a_i,-a_j)\),
has weight zero and positive degree.  Its strictly positive
valuation gives
\[
                       -b_i/a_i<-b_j/a_j.
 \tag{UCE.11}
\]
Thus the largest of the finitely many lower bounds is strictly
smaller than the smallest upper bound.  If one set is empty there
is only a one-sided bound.  Choose rational \(a\) in the resulting
open interval and choose \(t\) of valuation \(a\) in the given
algebraic closure.  Images equal to zero cause no constraint.
The direct grading in (UCE.9) makes (UCE.10) an actual algebra
map, and it has the required values.  This proves all cases,
including one-sided and zero images.
% END UCE14

% BEGIN UCE15
\paragraph{UCE15. The rescaled point extends continuously into one finite coefficient DVR.}
\label{claim:UCE15}
The map \(\xi_t\) of UCE14 extends to a local algebra map
\(f:S\to A\), where \(A\) is a subalgebra of the supplied
\(\overline{\mathbf Q}_\ell\), a complete DVR finite over
\(\mathcal O\).  Its inclusion into the algebraic closure agrees
with \(\xi_t\) on \emph{every} element of \(E_S\).

\emph{Proof.}
Choose a finite extension \(L/\operatorname{Frac}\mathcal O\)
inside the supplied algebraic closure containing \(t\), the
finitely many \(\xi(z_i)\), and the images of a finite set of
\(\mathcal O\)-module generators of \(S_0\).  By UCE13 all
of \(\xi_t(E_S)\) lies in \(L\).  Its generators over
\(\xi(S_0)\) have positive valuation, and \(\xi(S_0)\)
is integral, so this image lies in the ring of integers \(A\)
of \(L\).  This is a complete DVR finite over \(\mathcal O\)
and contains its original image.

Put \(b_i=\xi_t(z_i)\in\mathfrak m_A\).  Substitution
\(X_i\mapsto b_i\) gives convergent evaluation on the full
power-series ring \(P\) in (UCE.5).  To see that it factors
through \(S\), first take a homogeneous series \(F_j\).
Its truncations converge to its image in \(S_j\).  By UCE13
this is convergence in a finite \(\mathcal O\)-module, and
the \(\mathcal O\)-linear map \(\xi_t|_{S_j}\) is continuous.
Evaluation of \(F_j\) therefore agrees with that map, not only
with its values on polynomial truncations.
If \(F\in\ker(P\to S)\), every homogeneous projection
\(F_j\) is in the kernel by UCE6, so evaluates to zero.
The evaluation of \(F\) is the sum of these evaluations:
the variables all have norms bounded by one fixed \(q<1\),
so terms of total degree at least \(m\) have norm at most
\(q^m\), independently of their weight.  Grouping the convergent
sum by weights is thus legitimate.  Hence \(F\) evaluates to
zero and evaluation gives an actual map \(f:S\to A\).

The maximal ideal of \(S\) is generated by \(\pi,z_1,\ldots,z_d\),
and all their images lie in \(\mathfrak m_A\).  Thus \(f\)
is local and continuous.  The homogeneous argument above gives
the asserted equality on every \(S_j\), and (UCE.9) gives it
on the entire original entry algebra.  This establishes finite
coefficient realization without assigning an unjustified
continuity assumption to the initial map \(\xi\).
% END UCE15

% BEGIN UCE16
\paragraph{UCE16. The full original ordinary representative of every entry point.}
\label{claim:UCE16}
Assume \(\mathrm{FinA}(D,g_1)\).  For every original \(U,S\),
every map \(\xi:E_S\to\overline{\mathbf Q}_\ell\), and every
actual \(r:G\to\mathcal G_3(\overline{\mathbf Q}_\ell)\)
with \(\mathrm{IsPush}(\xi,r_S,r)\), there are a finite
coefficient DVR \(A\) inside the same algebraic closure, a unit
\(t\in\overline{\mathbf Q}_\ell^\times\), and an actual
ordinary point \(r':G\to\mathcal G_3(A)\) such that
\[
                    r'_\mathrm{incl}=T(t)\,r\,T(t)^{-1}
                        \quad\hbox{in }\mathcal G_3.
 \tag{UCE.12}
\]
This is the full literal \texttt{exists\_ordPt\_of\_point}.

\emph{Proof.}
Use UCE14 and UCE15, and set \(r'=f_*r_S\).  UCE5 proves
all its ordinary-point fields with the same marked \(g_1\)
and the same valuation ring \(B\).  Every entry of weight
\(w_i-w_j\) is multiplied by \(t^{w_i-w_j}\) under
\(\xi_t\).  These are exactly the entries of the conjugate
in (UCE.12).  Its component character is unchanged by
\(\mathrm{IsPush}\).  Finally \(\theta(T(t))=T(t)\), so
the same entry formula holds on the disconnected component as
well as on \(H\).  Thus (UCE.12) is equality of the actual
full group homomorphisms, not merely an isomorphism of their
semisimplifications or a trace comparison.
% END UCE16

\subsection{The exact automorphy consumers on the common marked datum}

% BEGIN UCE17
\paragraph{UCE17. Every irreducible ordinary point stays irreducible on every open subgroup.}
\label{claim:UCE17}
Let \(A\subset\overline{\mathbf Q}_\ell\) be an original finite
coefficient DVR and let \(r'\) be an \(\mathrm{OrdPt}\) for
\(D,g_1,B\).  If its restriction to \(H\) is irreducible,
then its restriction to every open subgroup \(V\subset H\)
is irreducible.  This is precisely \texttt{matIrr\_open\_gd}.

\emph{Proof.}
Replace \(V\) by its normal core \(V_0\) in \(G\); it
has finite index since \(G\) is profinite.  Choose an inertia
element with cyclotomic value \(1+\ell\), using total ramification
of every \(\ell\)-power cyclotomic extension and compactness.
Raise it to a positive multiple of \([G:V_0]\) and of the
order of the finite character part.  The resulting
\(u\in I_B\cap V_0\) has \(\Psi(u)\) of infinite order.
The three inertia eigenvalues \(\Psi(u),1,\Psi(u)^{-1}\)
are therefore distinct.

For clarity, the required finite-index Clifford argument is
elementary.  An irreducible \(H\)-module restricted to the normal
subgroup \(V_0\) is a sum of simple modules: take a minimal
nonzero \(V_0\)-submodule and sum its finitely many \(H\)-translates.
The sum is the whole module by irreducibility, and a sum of simple
submodules is semisimple.  Its isotypic components are permuted
transitively by \(H\).  Since the dimension is the prime three,
a reducible restriction is either scalar, or the sum of three
distinct character lines permuted transitively.  The distinct
eigenvalues of \(u\) rule out the scalar case.

Diagonalize \(r'(u)\) inside the ordered upper triangular
ordinary basis.  The diagonalizing change can be upper triangular,
by solving its three eigenvector equations successively.  The
three lines just found are its three eigenlines.  Every element
of \(H\) acts by a monomial matrix in this basis.
Take UCE1's \(x\in I_B\) with
\(\overline\Psi(x)^2\ne1\).  Its matrix is also upper
triangular, hence its permutation of the three lines is the
identity.  A transitive permutation group on three letters contains
a three-cycle.  Choose \(g\in H\) inducing one.  Thus
\[
                        \operatorname{tr}r'(gx^m)=0
                                      \quad(m\ge0).
 \tag{UCE.13}
\]
Reduce this equality in the original coefficient DVR.  With
\(a=\overline\Psi(g)\) and \(b=\overline\Psi(x)\), it becomes
\(ab^m+1+(ab^m)^{-1}=0\), or
\(a^2b^{2m}+ab^m+1=0\).  The equations for \(m=0,1\)
give \(a(b+1)+1=0\), since \(a(b-1)\ne0\).
Those for \(m=0,2\) give \(a(b^2+1)+1=0\), since
\(a(b^2-1)\ne0\).  Subtraction gives
\(ab(b-1)=0\), a contradiction.  Therefore the restriction
to \(V_0\), and hence to \(V\), is irreducible.
No additional small-prime restriction or cyclotomic-kernel
replacement was used.
% END UCE17

% BEGIN UCE18
\paragraph{UCE18. Both automorphic conclusions return along the actual torus conjugacy.}
\label{claim:UCE18}
Suppose (UCE.12) holds.  Type I is preserved, and an
\(\iota\)-automorphic conjugate-self-dual conclusion for
\(\rho_K(r')\epsilon^{-1}\) gives the same conclusion for
\(\rho_K(r)\epsilon^{-1}\).  Type II and its middle-entry-one
condition are preserved as well.  For every original extension
\(W\) of the outer corner of \(r\), set
\[
 D(t)=\operatorname{diag}(t,t^{-1}),\qquad
                         W'=D(t)WD(t)^{-1}.
 \tag{UCE.14}
\]
Then \(W'\) is the corresponding original extension of the
outer corner of \(r'\).  If it is attached, after the required
\(\epsilon^{-e}\) twist, to a newform of weight \(2e+1\),
the same newform and the same \(\iota\) attach to \(W\).

\emph{Proof.}
On \(H\), (UCE.12) is ordinary matrix conjugacy by \(T(t)\).
It identifies invariant subspaces and the actual representations
\(\rho_K\), and commutes with their scalar cyclotomic twists.
The defining automorphic datum is transported without change:
keep its cuspidal representation, regular algebraic infinity type
and conjugate-self-duality; transport its invariant complements
to preserve semisimplicity; and use invariance of characteristic
polynomials to preserve every member of the specified cofinite
Frobenius attachment set.  These are all the original
\(\mathrm{IsAutCsd}\) fields.

For Type II, all four odd-weight zero entries remain zero, and
the outer block is conjugated by \(D(t)\).  The middle entry
has weight zero and remains one.  The outer-corner equality in
the definition of \(\mathrm{IsExt}\) is thus satisfied by
\(W'\) at every element of \(H\), and its trace at the
specified complex conjugation is still zero by invariance of trace.
Conjugacy by \(D(t)\) is also an isomorphism of the actual full
two-dimensional group homomorphisms.  It commutes with
\(\epsilon^{-e}\).  Transporting invariant complements,
trivial inertia and every Frobenius characteristic polynomial
in the newform attachment preserves that attachment at the same
\(\iota\).  This argument applies to every supplied \(W\);
it does not choose a preferred extension or omit its other possible
quadratic twist.
% END UCE18

% BEGIN UCE19
\paragraph{UCE19. The original joint package from one common ordinary end datum.}
\label{claim:UCE19}
Suppose the original \(\mathrm{GdEnd}\) is supplied.  Then the
original \(\mathrm{AutomorphyUnramifiedCoefficients}\) holds.
More precisely, for a fixed original \(D\), fixed \(\iota\),
and actual local valuation ring \(B\), a single marked \(g_1\)
with the three original fields
\[
        \mathrm{FinEnd},\qquad
        \mathrm{AutEnd3}(K,\iota),\qquad
        \mathrm{AutEnd2}(e,\iota)
 \tag{UCE.15}
\]
supplies \(\mathrm{FinA}\), \(\mathrm{AutB}\), and
\(\mathrm{LemW}\) together, with exactly the original quantifiers.

\emph{Proof.}
Define \(B\) by the supplied local embedding of \(\Omega\).
It lies over \(\ell\).  Invoke \(\mathrm{GdEnd}\) once, with
the actual \(D\), \(B\), and fixed \(\iota\), obtaining the
single \(g_1\) in (UCE.15).  UCE12 supplies
\(\mathrm{FinA}(D,g_1)\) for all \(U,S\).

For \(\mathrm{AutB}\), now take any original \(U,S\), any
entry-algebra map \(\xi\), any actual pushed representation
\(r\) with its prescribed multiplier, and its Type-I hypothesis.
UCE16 gives an actual finite-coefficient ordinary point \(r'\)
conjugate to it by \(T(t)\).  UCE18 preserves Type I, and
UCE17 supplies irreducibility on every open subgroup of
\(H\), exactly the stronger antecedent in
\(\mathrm{AutEnd3}\).  Apply that field to the actual
\(\rho_K(r')\) obtained by restriction through the given
embedding \(K\hookrightarrow\Omega\).  It returns the original
\(\mathrm{IsAutCsd}\) assertion for
\(\rho_K(r')\epsilon^{-1}\).  UCE18 returns it to
\(\rho_K(r)\epsilon^{-1}\), which is precisely
\(D.\mathrm{IsAutI}(\iota,r)\).

For \(\mathrm{LemW}\), take the full original Type-II data,
including the middle-entry-one hypothesis and an arbitrary
\(W\) satisfying \(D.\mathrm{IsExt}(r,W)\).  Apply UCE16
to the same entry point.  UCE18 gives the Type-II ordinary
point with the same middle entry, and the actual full extension
\(W'\) in (UCE.14).  These are every hypothesis of
\(\mathrm{AutEnd2}\), including the trace-zero condition at
the original \(\sigma\).  That field supplies a newform of
weight \(2e+1\) attached to \(W'\epsilon^{-e}\).
UCE18 gives its attachment to \(W\epsilon^{-e}\), the literal
\(D.\mathrm{IsModular}(\iota,W)\).  Since \(W\) was arbitrary,
the universal extension quantifier in \(\mathrm{LemW}\) is retained.

Neither \(g_1\) nor \(\iota\) was changed in these arguments.
Only after their common choice were the local rings, normalized
rings, entry maps and pushed representations quantified.  Thus
the three fields hold simultaneously.  In particular they hold
when the supplied coefficient datum has \(e_{\mathcal O}=1\),
as required by the original joint target.  The implication uses
no extra ramification restriction on its intermediate finite DVRs.
% END UCE19

\subsection{Exact remaining producer and source boundary}

The closed results of this chapter are the local, ring, curve and
ordinary-representative assertions above, and their full original
same-object consumer.  A common \(g_1\) satisfying (UCE.15) has
not been constructed here.  Its three missing fields have different
roles: \(\mathrm{FinEnd}\) controls every characteristic-prime DVR
ordinary point using one finite-index subgroup; \(\mathrm{AutEnd3}\)
supplies actual polarized rank-three automorphy; and
\(\mathrm{AutEnd2}\) supplies actual newforms for every prescribed
outer-corner extension.  No result in this chapter obtains any of
those fields from finite weight zero or from an automorphic conclusion
that itself needs the joint package.

In the pinned original source the ordinary producer proceeds through
\texttt{GdFrameEx}, a common abelian CM frame, and
\texttt{LinkedEndE}, with its simultaneous rank-two and rank-three
runs, followed by the stated local and automorphic descent arguments.
The revised source route uses \texttt{ModUAlt}, the two actual seed
producers, and the corresponding frame and descent assertions.
The previously completed selected rank-four Hida support does not
assert the general \texttt{PfromAut} premise used to install that
\texttt{ModUAlt}.  Those remaining producers are recorded separately
in this component's contracts and source manifest.  Consequently
\texttt{AUT\_UNRAM\_COEF} remains open, and UCE19 is not recorded as
its whole-contract closure.
